\section{Problem statement}

Preparing academic documents remains a significant bottleneck for researchers and scholars.
\LaTeX\ offers unmatched quality and control but demands steep learning curves for formatting, structure, and syntax.
Writers must context-switch between content creation and markup management, dividing attention between what they want to say and how to say it in \LaTeX.
Even experienced users spend considerable time on repetitive tasks like organizing sections, placing and captioning figures, managing bibliography references, fixing formatting inconsistencies, and more.
This inefficiency cascades through the publishing pipeline as well.
Publishers can spend hours in back-and-forth corrections with authors over formatting errors, misaligned figures, and citation inconsistencies that could have been prevented with proper tooling.
Conversely, word processors like Microsoft Word or Google Docs offer ease of use but lack the precision, scalability, and professional output quality required for academic publishing, as they are built for a more general purpose audience.
This creates a painful gap: scholars and students either invest weeks mastering \LaTeX\ tooling before they can write effectively, or compromise on document quality by using less-capable alternatives.
Spartan Write eliminates this barrier by automating syntactic tasks through an AI agent, allowing first-time users and students to produce publication-quality documents without the traditional learning curve.

\subsection{Functional requirements}

\begin{itemize}
    \item Create \LaTeX\ projects from pre-defined templates.
    \item Open \LaTeX\ projects downloaded from other publishers.
    \item Converse with an AI agent that manages the document's source code.
    \item Allow direct source code manipulation if desired by the user.
    \item Support image uploads to the agent to attach to the document, or to parse and generate \LaTeX\ code, based on the user's request.
    \item Evaluate and configure the best LLM and provider for the task using a reproducible test suite.
    \item Deploy a user guide for grad-level students to quickly onboard and get started.
\end{itemize}

\subsection{Non-functional requirements}

\begin{itemize}
    \item Track AI usage and establish limits on usage using user accounts.
    \item Allow bring-your-own-key (BYOK) users to bypass established limits at their expense.
    \item Establish safeguards to ensure the agent remains within its scope.
    \item Design the agent to adhere to constraints of ethical use.
    \item Collect usage statistics for product iteration and insights.
    \item Split the application architecture in order to establish separation of concerns.
\end{itemize}

\subsection{Future improvements}

\begin{itemize}
    \item Support undo of agent's actions.
    \item Real-time collaboration for groups of users.
    \item Dynamic template repository.
    \item Integrate SyncTeX \cite{jlaurens_synctex} to improve document navigation by the agent.
\end{itemize}
